
\subsection{Contrato del servicio gRPC}

Para la comunicación mediante gRPC se definió el servicio \texttt{MessagingService} en el archivo \texttt{proto/messaging.proto}. Este contiene el método \texttt{SendMessage}, que recibe un mensaje de un nodo y devuelve un acuse de recepción.

El mensaje \texttt{Request} incluye \texttt{sender}, que identifica al emisor; \texttt{timestamp}, que lleva el valor del reloj de Lamport; y \texttt{message}, que contiene el texto enviado. La respuesta \texttt{Response} devuelve el estado, el nodo receptor y un detalle cuando ocurre un error.

\lstinputlisting[
  caption={Contrato del servicio gRPC},
  label={lst:messaging}
]{../proto/messaging.proto}

Las clases de los mensajes y las interfaces de comunicación se generaron con \texttt{grpcio-tools}, ejecutando desde la raíz del repositorio:

\begin{lstlisting}
python -m grpc_tools.protoc -I . --python_out=. --grpc_python_out=. proto/messaging.proto
\end{lstlisting}

Este comando genera \texttt{messaging\_pb2.py} y \texttt{messaging\_pb2\_grpc.py} dentro de \texttt{proto/}. El primero contiene los mensajes definidos y el segundo las clases que utilizan el cliente y el servidor.

